% GENERATED FILE. Edit canonical lessons, then run npm run generate:books.
% canonical-source-hash: fnv1a64:836e3c0f0a8142f1
% canonical-lessons: ES-C290-trigo, ES-C290-arroz, ES-C290-cebolla, ES-C290-tomate, ES-C290-miel

\chapter{What Is Eaten --- Wheat, Rice, Onion, Tomato, Honey}
\label{ch:lo-que-se-come}
\hlchaptermodality{290}

\begin{chapteropening}
\textbf{By the end of this chapter:} \emph{I can name wheat, rice, an onion, a tomato and honey in Spanish, and thank somebody for them.}

Complete ES-C290-miel: name all five foods in turn and thank somebody for each.
\end{chapteropening}

\section[el trigo]{el trigo — the grain that got rubbed}

\label{lesson:ES-C290-trigo}



\begin{quote}
Before a new run: what two Latin words hide in \emph{sombra}? (\emph{Sub} and \emph{umbra}.) And what was \emph{stella}? (A star.)
\end{quote}

\subsection*{What to know first}
\begin{itemize}
\raggedright
  \item el pan — what this becomes.
  \item el agua — the other half of what bread needs.
\end{itemize}

\begin{sounds}
\begin{itemize}
\raggedright
  \item \emph{TRI-go}, stress on the first part. One tap for the \emph{r}, and the \emph{g} is hard before \emph{o}. $\to$ reference
\end{itemize}
\end{sounds}

\begin{cousinweb}
\textbf{el trigo} — \textquotedblleft{}the wheat.\textquotedblright{}

Latin \textbf{triticum}, and that is built on \textbf{terere}, to rub or to wear away. Wheat was named for what you had to do to it: rub the grain out of the husk. The crop is named after the threshing, not the plant.

\emph{Terere} wore a path through English too. Something \textbf{trite} has been rubbed smooth by overuse. \textbf{Attrition} is wearing down. \textbf{Detritus} is what has been rubbed off. Even \textbf{detriment} is a wearing-away, though it has gone abstract and stopped admitting it.

A useful habit, this: when a Latin food word looks odd, ask what was done to the food. The answer is often the whole etymology.
\end{cousinweb}

\begin{grammarlens}[title={What you've built}]
One word, and a small piece of the ancient working day.

You have had \emph{el pan} for a long time. Now you have what it is made of, which means you can name a thing and the stuff it came from — a step up from naming things one at a time.
\end{grammarlens}

\subsection*{Your turn}
Say these aloud:

\begin{itemize}
\raggedright
  \item \emph{el trigo}
  \item \emph{el trigo}, then \emph{el pan} — the stuff, then the thing
  \item \emph{el pan de trigo}
\end{itemize}

\begin{tcolorbox}[breakable,colback=teal!4,colframe=teal!35!black,title={Before you move on}]
What did \emph{terere} mean? (To rub.) Which English words still rub? (\emph{Trite}, \emph{attrition}, \emph{detritus}.) And what is \emph{el trigo}? (The wheat.)
\end{tcolorbox}
\section[el arroz]{el arroz — a word that brought its own article}

\label{lesson:ES-C290-arroz}



\begin{quote}
Two back: what does \emph{la sombra} name? (The shade.) And what did \emph{terere} mean? (To rub.)
\end{quote}

\subsection*{What to know first}
\begin{itemize}
\raggedright
  \item el trigo — the grain before this one.
  \item el agua — what you cook it in.
\end{itemize}

\begin{sounds}
\begin{itemize}
\raggedright
  \item \emph{a-RROZ}, stress on the end. The \emph{rr} at the start of the second part is the rolled one, held. $\to$ reference
\end{itemize}
\end{sounds}

\begin{cousinweb}
\textbf{el arroz} — \textquotedblleft{}the rice.\textquotedblright{}

This one did not come from Latin. Arabic \textbf{ar-ruzz} crossed into Spain with the crop itself, and Spanish took the phrase whole, article and all. \emph{Ar-} was already \textquotedblleft{}the.\textquotedblright{} So \emph{el arroz} says it twice, in two languages, and nobody minds.

Once you know to watch for it, that stuck-on \emph{a-} is everywhere in the Arabic layer of Spanish: \emph{azúcar}, \emph{aceite}, \emph{almohada}, \emph{aduana}. Each one is still wearing an Arabic \emph{al-} that Spanish stopped being able to see.

English got the same grain by a different road — Greek \textbf{oryza}, then Latin, then Italian, then French, arriving as \textbf{rice}. Same plant, two routes, and the words no longer look like relatives.
\end{cousinweb}

\begin{grammarlens}[title={What you've built}]
Two grains, and two completely different histories sitting side by side.

\emph{El trigo} came down from Rome. \emph{El arroz} came in from the south with farmers. That is Spanish in miniature: a Latin skeleton with a large Arabic vocabulary laid over the top, and the joins still visible when you know where to look.
\end{grammarlens}

\subsection*{Your turn}
Say these aloud:

\begin{itemize}
\raggedright
  \item \emph{el arroz}
  \item \emph{el trigo}, then \emph{el arroz} — Rome, then the south
  \item \emph{arroz}, \emph{azúcar}, \emph{aceite} — three words carrying an Arabic \emph{al-}
\end{itemize}

\begin{tcolorbox}[breakable,colback=teal!4,colframe=teal!35!black,title={Before you move on}]
What does the \emph{ar-} in \emph{arroz} already mean? (The.) Name two more Spanish words carrying an Arabic article. (\emph{Azúcar}, \emph{aceite}, \emph{almohada}.) And which English word is the same grain? (\emph{Rice}.)
\end{tcolorbox}
\section[la cebolla]{la cebolla — a little something, permanently}

\label{lesson:ES-C290-cebolla}



\begin{quote}
Before the new one: what does the \emph{ar-} of \emph{arroz} already mean? (The.) And what was \emph{triticum} named for? (The rubbing, the threshing.)
\end{quote}

\subsection*{What to know first}
\begin{itemize}
\raggedright
  \item comer — what you eventually do with it.
  \item el arroz — what it usually goes into.
\end{itemize}

\begin{sounds}
\begin{itemize}
\raggedright
  \item \emph{ce-BO-lla}, stress in the middle. The \emph{c} before \emph{e} is soft, and the \emph{ll} is close to an English \emph{y}. $\to$ reference
\end{itemize}
\end{sounds}

\begin{cousinweb}
\textbf{la cebolla} — \textquotedblleft{}the onion.\textquotedblright{}

Latin had \textbf{caepa} for an onion, and \textbf{caepulla} for a small one. Spanish kept the small one and quietly dropped the full-sized word, so the language now calls every onion a little onion, however large.

That happens more than you would expect. A word with a diminutive on it gets used so often that the ending stops meaning anything, and the plain form falls away. You will meet the same \emph{-illa} and \emph{-illo} again shortly, still doing the job in words where it has not yet worn out.

English kept the other half. \textbf{Chive} came through French \emph{cive} from that same \textbf{caepa}, and a chive really is a small one — so between the two languages the pair got split down the middle.
\end{cousinweb}

\begin{grammarlens}[title={What you've built}]
Three things you can put in a pot, and a piece of machinery worth more than the three.

\emph{El trigo}, \emph{el arroz}, \emph{la cebolla}. Hearing a small ending as \textquotedblleft{}a little one of these\textquotedblright{} will unlock words you have not been taught, which is the whole point of taking them apart.
\end{grammarlens}

\subsection*{Your turn}
Say these aloud:

\begin{itemize}
\raggedright
  \item \emph{la cebolla}
  \item \emph{el arroz con cebolla}
  \item \emph{el trigo}, \emph{el arroz}, \emph{la cebolla}
\end{itemize}

\begin{tcolorbox}[breakable,colback=teal!4,colframe=teal!35!black,title={Before you move on}]
What was \emph{caepa}? (An onion.) What did \emph{caepulla} add? (Smallness.) Which English plant is the same root? (\emph{Chive}.)
\end{tcolorbox}
\section[el tomate]{el tomate — a word that came the other way}

\label{lesson:ES-C290-tomate}



\begin{quote}
Quick pass: what was \emph{caepa}? (An onion.) What does \emph{ar-} mean in \emph{arroz}? (The.) And what is \emph{el trigo}? (The wheat.)
\end{quote}

\subsection*{What to know first}
\begin{itemize}
\raggedright
  \item la cebolla — what it is usually cooked with.
  \item rojo — what colour it goes.
\end{itemize}

\begin{sounds}
\begin{itemize}
\raggedright
  \item \emph{to-MA-te}, stress in the middle. Three clean vowels, none of them swallowed. $\to$ reference
\end{itemize}
\end{sounds}

\begin{cousinweb}
\textbf{el tomate} — \textquotedblleft{}the tomato.\textquotedblright{}

Take this one apart and nothing Latin falls out, because there is nothing Latin in it. \textbf{Tomatl} is Nahuatl, the language of the Mexica, and the plant and its name travelled to Spain together in the sixteenth century.

Spanish shaved the \emph{-tl} into \emph{-te}, which is what it did to the whole set: \textbf{chocolatl} became \emph{chocolate}, \textbf{ahuacatl} became \emph{aguacate}, \textbf{coyotl} became \emph{coyote}. Then English took them from Spanish, which is why an English speaker already owns three Nahuatl words without knowing it.

So the traffic ran both ways. Spanish handed \emph{sombrero} to English and took \emph{tomate} from Mexico, and both words now sit in the same everyday vocabulary with no label saying where they have been.
\end{cousinweb}

\begin{grammarlens}[title={What you've built}]
Four foods, and three separate continents behind them.

\emph{El trigo} is Roman. \emph{El arroz} is Arabic. \emph{El tomate} is American. One short set of words, and the history of who fed whom is sitting inside it.
\end{grammarlens}

\subsection*{Your turn}
Say these aloud:

\begin{itemize}
\raggedright
  \item \emph{el tomate}
  \item \emph{tomate}, \emph{chocolate}, \emph{aguacate} — the same ending three times
  \item \emph{el arroz con tomate y cebolla}
\end{itemize}

\begin{tcolorbox}[breakable,colback=teal!4,colframe=teal!35!black,title={Before you move on}]
What language is \emph{tomate} from? (Nahuatl.) What happened to the \emph{-tl} ending? (It became \emph{-te}.) Name two more words that did it. (\emph{Chocolate}, \emph{aguacate}, \emph{coyote}.)
\end{tcolorbox}
\section[la miel]{la miel — the sweetness inside three English words}

\label{lesson:ES-C290-miel}



\begin{quote}
Before the last of the five: what language gave us \emph{tomate}? (Nahuatl.) What was \emph{caepa}? (An onion.) And what does \emph{ar-} mean? (The.)
\end{quote}

\subsection*{What to know first}
\begin{itemize}
\raggedright
  \item el pan — what you put it on.
  \item la flor — where it starts.
\end{itemize}

\begin{sounds}
\begin{itemize}
\raggedright
  \item \emph{MIEL}, one beat. The \emph{ie} is a single slide, and the \emph{l} is bright, never the dark English one at the end of a word. $\to$ reference
\end{itemize}
\end{sounds}

\begin{cousinweb}
\textbf{la miel} — \textquotedblleft{}the honey.\textquotedblright{}

Latin \textbf{mel}, four letters shorter than most of what we have taken apart, and Spanish did little to it beyond breaking the \emph{e} into \emph{ie} — the same break you met in \emph{cielo}.

\emph{Mel} is sweeter in English than it looks. Something \textbf{mellifluous} flows with honey, which is what \emph{fluere} is doing on the end of it. \textbf{Molasses} came through Portuguese \emph{melaço}, the honey-like leavings of sugar. And \textbf{marmalade} carries the Greek form \textbf{meli}, stuck to \emph{melon}, quince — a quince-honey, which is exactly what marmalade is.

There is a lesson hidden in that last one. A word can be perfectly transparent in its own language and completely opaque once it has crossed two borders.
\end{cousinweb}

\begin{grammarlens}[title={What you've built}]
Five: \emph{el trigo}, \emph{el arroz}, \emph{la cebolla}, \emph{el tomate}, \emph{la miel}.

Set out together they are a shopping list, and they are also four separate histories — Rome, Arabia, Mexico, and Rome again. Every one of the five arrived with an English relative attached, which is the quiet argument this whole approach rests on: you know more of this language than you have been taught.
\end{grammarlens}

\subsection*{Your turn}
Say these aloud:

\begin{itemize}
\raggedright
  \item \emph{la miel}
  \item \emph{el pan con miel}
  \item all five — \emph{el trigo}, \emph{el arroz}, \emph{la cebolla}, \emph{el tomate}, \emph{la miel}
\end{itemize}

\begin{tcolorbox}[breakable,colback=teal!4,colframe=teal!35!black,title={Before you move on}]
What was \emph{mel}? (Honey.) What does \emph{mellifluous} say, taken apart? (Flowing with honey.) And which breakfast preserve is a quince-honey? (\emph{Marmalade}.)
\end{tcolorbox}
